\chapter{Hybrid Models and Universal Differential Equations}
\label{ch:hybrid-ude}

\epigraph{\itshape The differential equation is the model. The neural
network is a closure.}{Rackauckas et al., SciML.ai}

\section{Universal differential equations}
A \emph{universal differential equation} (UDE) augments a known model
with a neural surrogate for the unknown
parts~\cite{rackauckas2020ude}:
\begin{equation}
\dot{\bm{x}} = \bm{f}(\bm{x},t;\mu) + \bm{N}_\theta(\bm{x},t),
\end{equation}
where $\bm{f}$ encodes physics we believe (kinematics, conservation),
and $\bm{N}_\theta$ absorbs missing dynamics --- closure terms,
unmodelled couplings, anomalous behaviour. Training is end-to-end via
adjoint sensitivity through ODE solvers
(\cref{ch:nodes,ch:numerics}).

\subsection{Why hybrid beats both extremes}
\begin{itemize}
\item Pure ODEs over-fit assumed structure and ignore observed
discrepancies.
\item Pure neural networks waste capacity rediscovering known physics
and generalise poorly out of distribution.
\item UDEs combine: low-data efficiency from physics, flexibility from
the network. The learned residual is often \emph{symbolic-regressible}
into a closed-form correction.
\end{itemize}

\subsection{Discovery loop}
\begin{recipe}{The Rackauckas--Edelman recipe}
\begin{enumerate}
\item Write down the known structural ODE/PDE.
\item Add $\bm{N}_\theta$ to the residual or specific terms.
\item Train end-to-end with adjoint gradients on observed data.
\item Apply SINDy / symbolic regression to $\bm{N}_\theta$ to extract
analytic closure.
\item Substitute back, re-train, validate.
\end{enumerate}
\end{recipe}

\section{Closure modelling}
Turbulence, plasma, and climate codes often suffer from unresolved
small scales whose effects must be parameterised. Neural closure models
replace classical Reynolds/Boussinesq closures with neural networks:
\begin{itemize}
\item \textbf{LES--ML closures}: subgrid-stress predicted by CNN/GNN.
\item \textbf{Climate parameterisations}: convection, cloud, radiation
schemes (e.g.\ NeuralGCM~\cite{kochkov2024neuralgcm}).
\item \textbf{Plasma turbulence}: gyrokinetic surrogates.
\end{itemize}

\paragraph{NeuralGCM.}
A 2024 hybrid global atmosphere model that couples a differentiable
dynamical core with neural parameterisations and matches state-of-the-art
weather forecast skill at a fraction of the compute~\cite{kochkov2024neuralgcm}.
A landmark for hybrid SciML at scale.

\section{Differentiable simulators}
A differentiable simulator exposes $\partial \text{output}/\partial
\text{input}$ through the entire physical pipeline (PDE solve, contact,
collision, etc.). Combined with neural augmentation, this gives
gradient-based parameter inference, control, and design optimisation:
\begin{itemize}
\item JAX-MD, JAX-CFD, JAX-Fluids, Phiflow, Modulus, Brax.
\item DiffTaichi, Warp, Nimble, Drake (robotics).
\item Diffrax + Equinox (Julia/JAX) for differentiable solvers.
\end{itemize}

\section{Kolmogorov--Arnold Networks (KAN)}
KANs~\cite{liu2024kan} replace the affine $\bm{W}\bm{x}+\bm{b}$ in MLPs
with learnable univariate spline activations on edges, motivated by the
Kolmogorov--Arnold representation theorem. They claim:
\begin{itemize}
\item Improved interpretability via per-edge function visualisation.
\item Sometimes superior accuracy on PDE / symbolic-regression
benchmarks at small parameter counts.
\item Direct symbolic extraction by fitting splines to elementary
functions.
\end{itemize}
\begin{pitfall}{KAN hype, KAN reality}
Independent benchmarks have shown that the gap between KANs and well-tuned MLPs is much smaller, and training cost is higher. Treat KANs
as one more architectural option, not a panacea.
\end{pitfall}

\section{Operator + physics hybrids}
\begin{itemize}
\item \textbf{PINO}: physics-informed neural operator combines FNO
training data with PDE residuals.
\item \textbf{PI-DeepONet}: residual loss for the trunk--branch
decomposition.
\item \textbf{Differentiable solver $\circ$ NN}: replace expensive
solver substeps with neural surrogates inside an otherwise classical
solver (multigrid smoothers, preconditioners, time-stepping).
\end{itemize}

\section{Equation-discovery pipelines}
Modern SciML pipelines string together:
\[
\text{data} \to \text{denoising / smoothing} \to \text{candidate libraries}
\to \text{sparse / symbolic fit} \to \text{validation}
\]
with neural surrogates entering wherever the bottleneck lies.
Auto-MoDe and AI-Hilbert (NeurIPS 2024) automate large parts of this
loop.

\section{Summary}
Hybrid SciML is arguably the most pragmatic part of the field:
mechanistic backbone $+$ data-driven correction. It dominates
applications where physics is partially known and computational budgets
are real --- climate, biology, plasma, materials.
